% Generated draft: requires mathematical and source review.
\documentclass[11pt]{article}
\usepackage[margin=25mm]{geometry}
\usepackage{fontspec}
\setmainfont{Latin Modern Roman}
\usepackage{amsmath,amssymb,mathtools,unicode-math}
\setmathfont{Latin Modern Math}
\usepackage{enumitem,needspace,xurl,hyperref}
\setlength{\parindent}{0pt}
\setlength{\parskip}{5pt}
\setlength{\emergencystretch}{3em}
\newcommand{\N}{\mathbb{N}}
\newcommand{\Z}{\mathbb{Z}}
\newcommand{\Q}{\mathbb{Q}}
\newcommand{\R}{\mathbb{R}}
\newcommand{\C}{\mathbb{C}}
\newcommand{\F}{\mathbb{F}}
\newcommand{\A}{\mathbb{A}}
\newcommand{\PP}{\mathbb{P}}
\newcommand{\E}{\mathbb{E}}
\newcommand{\eps}{\varepsilon}
\newcommand{\dd}{\,\mathrm{d}}
\begin{document}
\section*{UnsolvedMath ID: 30001355}
\textbf{Source problem code:} catalogue entry
\subsection{Definitions and mathematical statement}
For a group \(\Gamma\), define its commutator subgroup and its abelianisation by
\[
[\Gamma,\Gamma]
=\bigl\langle ghg^{-1}h^{-1}:g,h\in\Gamma\bigr\rangle,
\qquad
\Gamma^{\mathrm{ab}}=\Gamma/[\Gamma,\Gamma].
\]
The quotient \(\Gamma^{\mathrm{ab}}\) is an abelian group.

The supplied source asks:
\begin{quote}
What can be said about the abelianisation of arithmetic groups, particularly arithmetic groups of orthogonal type?
\end{quote}
Thus the requested task is to determine or describe \(\Gamma^{\mathrm{ab}}\) for arithmetic groups \(\Gamma\), with particular emphasis on groups of orthogonal type. This is a broad research question rather than a specified conjectural assertion. The supplied record does not fix the underlying algebraic group, base field, ring of integers, quadratic form or lattice, dimension, signature, or arithmetic subgroup. It also does not specify whether full orthogonal groups, special orthogonal groups, or particular subgroups are intended. Consequently, no particular domain, family, or additional hypotheses are imposed here, and no more specific universally quantified classification statement can be recovered from this record alone.
\subsection{Short English statement}
The question asks for the structure of the quotient of an arithmetic group by its commutator subgroup, especially for arithmetic groups of orthogonal type. The intended family of groups and the precise kind of description sought remain unspecified.
\subsection{Sources}
\begin{itemize}
\item Supplied local UnsolvedMath registry record, permanent ID \(30001355\), title \emph{Abelianizations of Arithmetic Orthogonal Groups}, exact source code \texttt{catalogue entry}, in the collection \emph{Oberwolfach Reports Open Problems}. The proposed year is recorded as 2009; no individual proposer is supplied.
\item Supplied report citation: \emph{Complex Algebraic Geometry (2010)}, Oberwolfach Reports, pp.\ 2483--2544, DOI 10.4171/owr/2009/44. \url{https://doi.org/10.4171/owr/2009/44}. The background describes Oberwolfach Reports as documenting mathematical workshops and their open-problem discussions.
\item Supplied literature reference: Gritsenko--Hulek--Sankaran, ``orthogonal groups.'' \url{https://arxiv.org/abs/0810.1614}. The supplied assessment attributes to these authors computations of the commutator index for many indefinite orthogonal groups and applications to modular varieties; the record does not state the precise hypotheses or results.
\item Supplied literature assessment dated exactly \texttt{2026-08-21}: the record labels the problem \texttt{partially\_solved} and describes the broad extracted abelianisation question as not fully resolved. This is supplied dated curation derived from the OWR collection, not an independently verified literature assessment. Its listed sources are the Gritsenko--Hulek--Sankaran reference and the original OWR report linked above; time-sensitive status and cited claims require independent verification.
\item The registry labels the statement \texttt{exact} and reports that its mathematical content was checked against the official Oberwolfach report, with notation, context, and question form normalized. This verification claim is supplied metadata and has not been independently checked here.
\end{itemize}
\end{document}
